% !PRINTED 534
\weberpage{534}
%<<CONT
we further denote two opposite classes by $(A, \pm 2B, C)$, then the
corresponding values of $f(\omega)$ are conjugate imaginary, their product
therefore positive, and it follows accordingly from these determinations
that the product
\[
\Pi f(\omega)
\]
has a positive real value which is a power of 2. Denoting the class number
by $h$, we set this power $= 2^{h\tau}$, so that
\weq{1}{\Pi\,\frac{f(\omega)}{2^\tau} = 1.}

It will be our task to determine this exponent $\tau$. But we first premise
the following remark, which lends this task an increased interest.

In consequence of the equation [§ 54, (8)]:
\weq{2}{f(\omega)^{24} - \gamma_2(\omega) f(\omega)^8 - 16 = 0}
it follows that, if $\omega$ is the root of a form of the discriminant
$-4m$ belonging to the class $k$, $f(\omega)$ is an \gesperrt{integral algebraic
number}; and since $f(\omega)^8$ in (2) may also be replaced by
$-f_1(\omega)^8$, $-f_2(\omega)^8$, the quantities $f_1(\omega)$,
$f_2(\omega)$ are likewise integral algebraic numbers. Consequently so is
\[
\frac{\sqrt{2}}{f(\omega)} = f_1(\omega) f_2(\omega).
\]

If $p$ is an odd prime of which $-m$ is a quadratic residue, and if $p$ is
representable by the forms of the class $l$ (of the discriminant $-4m$),
then, upon a suitable determination of $c$ according to § 118,
\[
\frac{c+\omega}{p}
\]
is the root of a form belonging to the compounded class $lk$, and (unless
it happens that $p=3$) $c$ may be assumed divisible by 48. If we therefore
set:
\[
u = f(\omega), \quad v = f\left(\frac{c+\omega}{p}\right),
\]
then $uv$ as well as $2:uv$ is an \gesperrt{integral algebraic number}.

If therefore we set, following § 74,
\[
\begin{aligned}
B &= (uv)^s + \left(\frac{2}{p}\right)\frac{2^s}{(uv)^s},\\
A &= \left(\frac{u}{v}\right)^r + \left(\frac{v}{u}\right)^r
\end{aligned}
\]
%CONT>>
